\documentclass[11pt,a4paper]{article}
\usepackage[T1]{fontenc}
\usepackage{lmodern}
\usepackage{amsmath,amssymb,amsthm,mathtools}
\usepackage[margin=26mm]{geometry}
\usepackage{microtype,booktabs,array,xcolor}
\usepackage[hidelinks]{hyperref}
\usepackage{fancyhdr,enumitem}
\definecolor{epink}{HTML}{452C43}
\hypersetup{pdftitle={Field-uniform nonintegrability of spin chains},pdfauthor={Anonymous},pdfsubject={Evidence Press research release candidate; version 1.0.1-candidate; 28 September 2026}}
\pagestyle{fancy}\fancyhf{}
\fancyhead[L]{\small\sffamily EVIDENCE PRESS\quad /\quad MATHEMATICAL PHYSICS}
\fancyhead[R]{\small\sffamily 28 September 2026}
\fancyfoot[L]{\small\sffamily Field-uniform nonintegrability\quad v1.0.1}
\fancyfoot[R]{\small\thepage}
\setlength{\headheight}{15pt}
\setlength{\parskip}{4pt}\setlength{\parindent}{0pt}
\setlength{\emergencystretch}{3em}
\numberwithin{equation}{section}
\newtheorem{theorem}{Theorem}[section]
\newtheorem{lemma}[theorem]{Lemma}
\newtheorem{proposition}[theorem]{Proposition}
\newtheorem{corollary}[theorem]{Corollary}
\theoremstyle{definition}\newtheorem{definition}[theorem]{Definition}
\theoremstyle{remark}\newtheorem{remark}[theorem]{Remark}
\DeclareMathOperator{\Tr}{Tr}\DeclareMathOperator{\rank}{rank}
\DeclareMathOperator{\diag}{diag}\DeclareMathOperator{\Disc}{Disc}
\DeclareMathOperator{\End}{End}\DeclareMathOperator{\len}{len}
\DeclareMathOperator{\im}{im}\DeclareMathOperator{\ad}{ad}
\newcommand{\spanC}{\operatorname{span}_{\mathbb C}}
\newcommand{\Q}{\mathbb Q}\newcommand{\C}{\mathbb C}\newcommand{\R}{\mathbb R}
\newcommand{\g}{\mathfrak{sl}_d(\C)}\newcommand{\id}{\operatorname{id}}
\newcommand{\Ob}{\mathcal O}\newcommand{\tens}{\otimes}
\newcommand{\Frob}{\mathrm F}\newcommand{\op}{\mathrm{op}}
\begin{document}
\thispagestyle{empty}
{\color{epink}\sffamily\small EVIDENCE PRESS\enspace /\enspace RESEARCH RELEASE CANDIDATE}
\vspace{7mm}

{\color{epink}\LARGE\bfseries Field-uniform nonintegrability\par of spin chains\par}
\vspace{4mm}
{\large An exact three-site quotient, symmetry averaging and a complete bilinear--biquadratic field-existence criterion\par}
\vspace{5mm}
{\small Anonymous\quad\textbullet\quad Version 1.0.1-candidate\quad\textbullet\quad 28 September 2026}
\vspace{3mm}

{\small DOI: \url{https://doi.org/10.5281/zenodo.23021227}\par}
\begin{abstract}
We study local conservation for a fixed symmetric nearest-neighbour interaction while allowing arbitrary, site-dependent Hermitian on-site fields. Under two leading-term rigidity conditions, automatic at full traceless interaction rank, we construct an exact finite-dimensional quotient that eliminates every field and every two-site correction from the three-site commutator equations. Nonvanishing of its obstruction excludes conserved quantities of every exact contiguous range $3\le k\le\lfloor N/2\rfloor$, for every periodic $N\ge6$ and every choice of fields. Conversely, vanishing is equivalent to the existence of a uniform-field three-site quantity whose commutator has range at most two. Thus the quotient is complete for this almost-conservation test, not for integrability. For qutrits it has only $36$ real correction and field variables.

A multiplicity-free symmetry theorem evaluates the quotient without solving those equations. It shows that arbitrary fields cannot remove the full three-site obstruction of an invariant interaction. Applied to the spin-1 bilinear--biquadratic interaction $uB+v(B^2-4I/3)$, the obstruction is $u(u^2-v^2)W$, with $\|W\|_F^2=40$. Exact rank certificates close both rank-deficient lines. Consequently fields can admit a higher-range local charge for a nonzero interaction in this family if and only if $u(u^2-v^2)=0$; zero fields suffice on those three lines. Off them, the entire local commutant is energy, identity and the uniform spin generators commuting with every field. We also certify generic exclusion on both hypersurfaces missed by an earlier partial-trace obstruction, while retaining its unrestricted generic theorem. All finite certificates and separately implemented exact checks accompany the paper.
\end{abstract}

\textbf{Scope and assurance.} The all-range implication uses Hokkyo's revised reduction theorem. The paper does not classify arbitrary fields on the three integrable interaction lines, the complete general qutrit exceptional set, or quasilocal charges. The supplied referee report concerns version 0.1; the new theorems in this version have exact computational corroboration but have not received a fresh independent specialist review.

\section{Setting, questions and prior results}\label{sec:setting}

Let $F_1,\ldots,F_m$ be a Hermitian, Hilbert--Schmidt orthonormal basis of the traceless $d\times d$ matrices, where $m=d^2-1$ and $\Tr(F_aF_b)=\delta_{ab}$. Consider
\begin{equation}\label{eq:ham}
 H_N=\sum_{i=1}^N h_{i,i+1}+\sum_{i=1}^N f_i,
 \qquad h=\sum_{a,b=1}^m J_{ab}F_a\tens F_b,
 \qquad J=J^T\in M_m(\R).
\end{equation}
The chain is periodic, $N\ge6$, and every $f_i$ is an arbitrary traceless Hermitian matrix on site $i$. The same interaction $h$ occurs on every bond. In particular, $h$ has zero one-site partial traces and is invariant under the site swap. Fields may have unbounded norms, differ at every site and break all symmetries of $h$.

Every Hermitian swap-symmetric two-site density has a unique decomposition
\[
 \widetilde h=cI\tens I+h+\tfrac12(f\tens I+I\tens f),
 \qquad \Tr_1h=\Tr_2h=0.
\]
For $d=3$, the homogeneous-field class therefore has $36+8=44$ real parameters after an irrelevant scalar is removed. An orthogonal change of operator coordinates need not be a physical on-site unitary; none of our arguments equates these transformations.

\begin{definition}[Contiguous range]
An operator has range at most $k$ if it is a sum of operators supported on cyclic intervals of at most $k$ consecutive sites. Its exact range, denoted $\len(Q)$, is the smallest such $k$. Coefficients may depend on the site. A conserved quantity satisfies $[H_N,Q]=0$. A string $A_iI_{i+1}B_{i+2}$ with nonzero traceless endpoints has range three, not two.
\end{definition}

We ask a question with an explicit quantifier over the fields: \emph{can any choice of the $f_i$ admit a local conserved quantity of range at least three?} A negative answer is stronger than failure of one current to commute at zero field. Our tests must exclude all shorter-range corrections and charges without translation symmetry.

\subsection{Established ingredients and the present contribution}
The local-conservation testing programme goes back to Grabowski and Mathieu~\cite{GrabowskiMathieu}. Hokkyo's revised theorem reduces an all-range exclusion to the absence of an exact-range-three $Q$ satisfying $\len[H_N,Q]\le2$, provided its two hypotheses hold~\cite{HokkyoTest}. Here almost-conservation means that the commutator has support of range at most two; it does not mean a small commutator norm. Excluding this almost-conserved quantity is stronger than excluding an exactly conserved three-site charge. We prove the hypotheses at full interaction rank and certify them on the two additional rank-deficient lines needed below.

Earlier spin-1 work establishes the field-free bilinear--biquadratic classification and the classification with a uniform uniaxial on-site anisotropy, as well as generic anisotropic models with additional internal symmetries~\cite{ParkLee,HSC}. Further $U(1)$-symmetric families with vanishing quadrupolar couplings are treated by Sengoku and Watanabe~\cite{SW}. Shiraishi and Yamaguchi prove an isotropic arbitrary-spin dichotomy~\cite{SY}. Those classifications are important comparators: our isotropic result does not claim to discover the zero-field integrable lines. Its additional content is the quantifier over \emph{every site-dependent Hermitian on-site field}, the exact residual commutant, and a symmetry principle that gives this conclusion.

\paragraph{The precise extension of the isotropic predecessor.}
Shiraishi and Yamaguchi~\cite{SY} treat a translation-invariant, zero-field isotropic Hamiltonian. Their proof for a nonzero three-support obstruction already constrains the two-site correction using $SU(2)$ and time-reversal symmetry. Thus removing the correction by symmetry is not, by itself, a new idea. Here the Hamiltonian may break both internal and translation symmetries through arbitrary site-dependent fields. We average the \emph{affine local almost-conservation equations}, including their field variables, rather than assert that the conserved operator inherits a symmetry of this Hamiltonian. Under the irreducible, multiplicity-free hypotheses, averaging kills both the swap-odd correction and the traceless field while retaining the obstruction. This is the step that upgrades a zero-field test to an exclusion uniform over all fields. The quotient-consistency theorem and the residual-commutant theorem are distinct further statements; neither asserts complete integrability of a general quotient-zero interaction.

The other contribution is an exact quotient of the three-site equations. Translation averaging is applied to their leading commutator components, not to a putative conserved charge under a Hamiltonian whose fields break translation symmetry. Reversal then reduces the qutrit correction and field variables to $28+8=36$. The quotient detects concrete full-rank interactions, and relatively generic families, on which the earlier cubic partial trace is identically uninformative.

\subsection{Relation to the Reshetikhin condition}\label{sec:reshetikhin}
Put
\begin{equation}\label{eq:CD}
 C=[h_{12},h_{23}],\qquad D=[h_{12}+h_{23},C].
\end{equation}
The full Reshetikhin condition is the existence of a two-site $X$ with
\begin{equation}\label{eq:reshetikhin}
 D=X_{23}-X_{12}.
\end{equation}
Let $\pi(x)=x-\Tr(x)I/d$ and define
\begin{equation}\label{eq:delta}
 \Delta(h)=(\pi_1\pi_3)D,\qquad \Ob(h)=\Tr_2D.
\end{equation}
Partial traces are unnormalised. The operator $\Delta$ retains the full middle-site matrix space, including its scalar direction. Equation~\eqref{eq:reshetikhin} implies $\Delta=0$, and hence $\Ob=0$. In the last implication we use that $\Ob$ already has traceless endpoints, so $\Tr_2\Delta=\Ob$. The converse from $\Ob=0$ fails; the AKLT example below has $\Ob=0$ but $\|\Delta\|_F^2=2560/81$.

Hokkyo proves that the specific boost-generated third conservation law generates a commuting hierarchy~\cite{HokkyoLaw}. Sanatani, Shiraishi and Ishii prove a sufficiency result for a regular difference-form Yang--Baxter construction from the full Reshetikhin condition with a diagonalisability hypothesis~\cite{SSI}. These results do not identify vanishing of a partial projection with integrability, nor do they alone give our arbitrary-field exclusion. The current paper uses their framework for positioning; its sole external all-range proof dependency is~\cite{HokkyoTest}.

\section{Leading rigidity and the all-range reduction}\label{sec:rigidity}

We separate the interaction hypotheses from any choice of fields.

\begin{definition}[Leading-rigid interaction]\label{def:rigid}
A nonzero interaction $h\in\g\tens\g$ is leading-rigid when it has both properties:
\begin{enumerate}[label=(\roman*),leftmargin=7mm,itemsep=1pt]
\item $[x\tens I,h]=0$ or $[I\tens x,h]=0$, with $x\in\g$, implies $x=0$.
\item If $a,b\in\g\tens\g$ satisfy
\begin{equation}\label{eq:pairlocal}
 [a_{12},h_{23}]=[h_{12},b_{23}],
\end{equation}
then $a=b=ch$ for one scalar $c$.
\end{enumerate}
\end{definition}

This local pair condition is sufficient for, and stronger in form than, the global two-site hypothesis of the external theorem. It is convenient because a fixed finite certificate establishes it for every chain length.

\begin{lemma}[Full rank supplies leading rigidity]\label{lem:fullrank}
If $J$ is invertible, $h$ is leading-rigid.
\end{lemma}
\begin{proof}
The left and right tensor factors of $h$ span all traceless matrices, so condition (i) follows from the scalar commutant of $M_d(\C)$. For (ii), write $a=\sum M_{rs}F_r\tens F_s$ and $b=\sum N_{rs}F_r\tens F_s$. Comparing the two endpoint factors in~\eqref{eq:pairlocal} gives
\[
 \left[\sum_sM_{rs}F_s,\sum_tJ_{tu}F_t\right]
 =\left[\sum_sJ_{rs}F_s,\sum_tN_{tu}F_t\right].
\]
Both families determined by $J$ are bases. Thus this equation is $[Lx,y]=[x,Ry]$ for two endomorphisms of $\g$. The pair-rigidity lemma in Appendix~\ref{app:pair} gives $L=R=c\id$, whence $M=N=cJ$.
\end{proof}

The scalarity mechanism is established quasicentroid theory, not a new Lie-algebra theorem. For a centreless perfect Lie algebra the quasicentroid equals the centroid; the latter is scalar for a finite-dimensional complex simple Lie algebra. We attribute this background to Leger and Luks, as recalled explicitly by Andruszkiewicz, Brzezi\'nski and Radziszewski~\cite{LegerLuks,ABR}. Appendix~\ref{app:pair} retains a self-contained proof for $\mathfrak{sl}_d$. The application to the leading-charge equations is what is needed here.

\begin{proposition}[Specialised Hokkyo reduction]\label{prop:reduction}
Suppose $h$ is leading-rigid. For any fields in~\eqref{eq:ham}, existence of a conserved quantity of exact range $3\le k\le\lfloor N/2\rfloor$ implies the existence of an exact-range-three $Q$ with $\len[H_N,Q]\le2$. Every such $Q$ has the form
\begin{equation}\label{eq:Qform}
 Q=c C_N+A,\qquad
 C_N=\sum_i[h_{i,i+1},h_{i+1,i+2}],\qquad
 c\ne0,\quad \len(A)\le2.
\end{equation}
In particular $C_N$ has exact range three.
\end{proposition}
\begin{proof}[Verification of the external hypotheses]
Condition (i) is exactly the injectivity assumption (4) of~\cite{HokkyoTest}. A strictly two-local operator $\sum_i a_i$ has no range-three part in its commutator precisely when~\eqref{eq:pairlocal} holds with $(a,b)=(a_i,a_{i+1})$ at every interval. Condition (ii) forces $a_i=c_i h$ and $a_{i+1}=c_i h$; adjacent equations give a common scalar. Hence the source's space $\mathcal B_{\le}^{(2)}$ equals $\C H_N^{(2)}$, its assumption (6). Theorem 1 then gives the range-three reduction, and the $2\to3$ isomorphism of Theorem 2 fixes the leading part to $cC_N$ and guarantees it is nonzero. The setup of that source explicitly allows site-dependent on-site terms.
\end{proof}

$C_N$ is anti-Hermitian; $iC_N$ is the corresponding Hermitian observable. Multiplying a charge by this scalar does not change conservation or its range. No finite-chain numerical calculation replaces this implication. All finite computations below verify its hypotheses or the local obstruction.

\section{An exact quotient that eliminates all fields}\label{sec:quotient}

Let $\mathfrak h_d$ be the real vector space of traceless Hermitian $d\times d$ matrices. Write $S$ for the swap on two sites and $R$ for reversal of three sites. Define
\[
 \mathcal A_-=
 \{a\in\mathfrak h_d\tens_{\R}\mathfrak h_d:SaS=-a\},
 \qquad \dim_{\R}\mathcal A_-=\frac{m(m-1)}{2}.
\]
Since $h$ is swap-symmetric, $RCR=-C$ and $RDR=-D$. Both $D$ and $\Delta$ are Hermitian. Define a real linear map
\begin{equation}\label{eq:T}
 \begin{split}
 T_h:\mathcal A_-\oplus\mathfrak h_d&\longrightarrow
 \mathfrak h_d\tens_{\R}\operatorname{Herm}(\C^d)\tens_{\R}\mathfrak h_d,\\
 T_h(a,f)&=i\bigl([h_{12},a_{23}]+[h_{23},a_{12}]\bigr)
       +[f_1+f_2+f_3,C].
 \end{split}
\end{equation}
Every term on the right has traceless factors at all three sites and is Hermitian and reversal-odd. In particular its endpoint projection is itself.

\begin{definition}[Field-eliminated obstruction]
The obstruction is the class of $\Delta(h)$ in the quotient by $\im T_h$. Equivalently, with the real Hilbert--Schmidt inner product on Hermitian operators, use the unique orthogonal residual
\begin{equation}\label{eq:residual}
 \mathcal U(h)=\Delta(h)-\operatorname{proj}_{\im T_h}\Delta(h).
\end{equation}
The zero/nonzero test is intrinsic; an orthonormal basis is not needed to implement it by exact ranks.
\end{definition}

\begin{theorem}[Exact field elimination]\label{thm:quotient}
For a leading-rigid, Hermitian, swap-symmetric interaction $h$, the following are equivalent:
\begin{enumerate}[label=(\alph*),leftmargin=7mm,itemsep=2pt]
\item $\Delta(h)\in\im T_h$, or equivalently $\mathcal U(h)=0$.
\item For some periodic $N\ge6$ and some site-dependent Hermitian fields, an exact-range-three $Q$ satisfies $\len[H_N,Q]\le2$.
\item There are fixed homogeneous fields $f_i=f$ and a fixed two-site correction density giving such a $Q$ for every periodic $N\ge6$.
\end{enumerate}
Consequently, $\mathcal U(h)\ne0$ excludes every conserved quantity of exact range $3\le k\le\lfloor N/2\rfloor$, simultaneously for all $N\ge6$ and all on-site fields.
\end{theorem}

\begin{proof}
Assume (b). Divide~\eqref{eq:Qform} by its leading coefficient and take its anti-Hermitian part. Since $H_N$ is Hermitian and $C_N^\dagger=-C_N$, this preserves both the leading part $C_N$ and the condition on the commutator. We may therefore write
\[
 Q=C_N+i\sum_i (a_i)_{i,i+1}+A^{(1)},
 \qquad a_i\in\mathfrak h_d\tens_{\R}\mathfrak h_d,
 \qquad \len(A^{(1)})\le1.
\]
The four-site contributions to $[H_N^{(2)},C_N]$ cancel pairwise by the Jacobi identity and commutation of disjoint bonds. Its remaining contribution on interval $i,i+1,i+2$ is $D_i$. Taking the exact range-three part on that interval gives
\begin{equation}\label{eq:localeq}
 \Delta+[f_i\tens I\tens I+I\tens f_{i+1}\tens I+I\tens I\tens f_{i+2},C]
 +i\bigl([h_{12},(a_{i+1})_{23}]+[h_{23},(a_i)_{12}]\bigr)=0.
\end{equation}
The one-site correction cannot contribute. Average these equations after identifying their three-site coordinate spaces. With
\[
 \bar f=N^{-1}\sum_i f_i,\qquad \bar a=N^{-1}\sum_i a_i,
\]
we obtain the same equation with each field equal to $\bar f$ and each correction equal to $\bar a$. This is an average of \emph{equations}; no symmetry of the inhomogeneous Hamiltonian is assumed.

Apply the reversal-odd projection $X\mapsto(X-RXR)/2$. It preserves $\Delta$ and the uniform-field term. It replaces $\bar a$ by $a_-=(\bar a-S\bar aS)/2$. We obtain
\[
 \Delta+T_h(a_-,\bar f)=0,
\]
which proves (a).

Conversely, if $\Delta+T_h(a,f)=0$, choose uniform fields $f$ and
$Q=C_N+i\sum_i a_{i,i+1}$. The same local calculation makes every range-three component vanish, while the four-site terms cancel as above. Thus $\len[H_N,Q]\le2$ for every $N\ge6$. Its leading part is nonzero by Proposition~\ref{prop:reduction}. This proves (c); (c) implies (b). The final conclusion follows from that proposition.
\end{proof}

\begin{corollary}[Dual certificate]\label{cor:dual}
If a Hermitian three-site $Z$ satisfies
\[
 \Tr(ZT_h(a,f))=0\quad\hbox{for all }a,f,
 \qquad \Tr(Z\Delta(h))\ne0,
\]
then the all-field exclusion in Theorem~\ref{thm:quotient} holds. Conversely such a $Z$ exists whenever $\mathcal U(h)\ne0$; one may take $Z=\mathcal U(h)$.
\end{corollary}

For $d=3$, $T_h$ has $28+8=36$ real columns. Its natural ambient space has dimension $8\cdot9\cdot8=576$, and its image is further restricted by reversal. The certifier uses a redundant, transparent coordinate system: real row-major entries of the $27\times27$ matrix followed by imaginary entries, giving $1458$ rows. Redundancy changes neither rank nor consistency.

\textbf{Meaning of completeness.} Theorem~\ref{thm:quotient} is complete for the existential, field-eliminated \emph{almost-conservation} test in (b). When the quotient vanishes, the residual range-one and range-two equations may still prevent exact conservation. The theorem does not identify its zero set with the integrable locus.

\section{Multiplicity-free symmetry removes the correction space}\label{sec:symmetry}

\begin{theorem}[Symmetry evaluation of the quotient]\label{thm:symmetry}
Let a compact group $G$ act irreducibly and unitarily on $V\cong\C^d$, and suppose $V\tens V$ is a multiplicity-free $G$-representation. Let $h\in\g\tens\g$ be Hermitian, $G$-invariant and leading-rigid. Then
\begin{equation}\label{eq:symres}
 \mathcal U(h)=\Delta(h).
\end{equation}
In particular, if $\Delta(h)\ne0$, no choice of site-dependent Hermitian fields admits conserved quantities of the ranges stated in Theorem~\ref{thm:quotient}.
\end{theorem}
\begin{proof}
The two-site commutant of $G$ is commutative because the representation is multiplicity-free. The swap belongs to it. Thus every invariant two-site operator, including $h$, commutes with the swap.

Let $\mathbb E_r$ be Haar averaging of conjugation by $G$ on $V^{\tens r}$. Irreducibility implies $\mathbb E_1(f)=\Tr(f)I/d=0$ for traceless $f$. An invariant two-site operator commutes with the swap, so averaging a swap-odd $a$ gives $\mathbb E_2(a)=0$. Invariance of $h$ and $C$ yields
\[
 \mathbb E_3(T_h(a,f))=T_h(\mathbb E_2(a),\mathbb E_1(f))=0.
\]
The operator $\Delta(h)$ is invariant, since both partial-trace projection and the commutators are equivariant. Haar averaging is the orthogonal projection onto invariant operators. Therefore $\Delta(h)\perp\im T_h$, proving~\eqref{eq:symres}. Theorem~\ref{thm:quotient} supplies the conclusion.
\end{proof}

The fields need not retain the symmetry: the averaging above acts on the linear correction and field terms, not on $H_N$ as a conserved system. This is the point that extends the invariant zero-field test to arbitrary on-site potentials.

\begin{corollary}[Isotropic interactions at full rank]\label{cor:spinS}
For the irreducible spin-$S$ representation, $V_S\tens V_S=\bigoplus_{L=0}^{2S}V_L$ is multiplicity-free. Every full-traceless-rank isotropic interaction with $\Delta(h)\ne0$ therefore has the arbitrary-field exclusion property.
\end{corollary}

The full-rank condition can be replaced by Definition~\ref{def:rigid}. We verify the needed rank-deficient spin-1 cases next. No assertion about every rank-deficient higher-spin interaction is needed.

\section{The complete spin-1 bilinear--biquadratic field-existence criterion}\label{sec:blbq}

Let $S_x,S_y,S_z$ be the spin-1 matrices and
\begin{equation}\label{eq:blbq}
 B=\sum_{\alpha=x,y,z}S_\alpha\tens S_\alpha,
 \qquad h(u,v)=uB+v(B^2-\tfrac43I),\qquad (u,v)\in\R^2.
\end{equation}
The scalar subtraction gives zero one-site partial traces. In a Cartesian spin basis, put
\[
 P|ab\rangle=|ba\rangle,\qquad
 Q_0=|\Omega\rangle\langle\Omega|,\qquad |\Omega\rangle=\sum_{a=0}^2|aa\rangle.
\]
Then $B=P-Q_0$, $B^2=I+Q_0$, and
\begin{equation}\label{eq:cartesian}
 h(u,v)=uP+(v-u)Q_0-\tfrac v3I.
\end{equation}
Here $Q_0$ is a two-site contraction, not a conserved quantity.

\begin{proposition}[Exact cubic and interaction spectrum]\label{prop:isocubic}
For the family~\eqref{eq:blbq},
\begin{align}
 \Ob(h(u,v))&=0,\label{eq:isoblind}\\
 \Delta(h(u,v))&=u(u-v)(u+v)W,
 &W&=\Delta(h(1,0)),& \|W\|_F^2&=40.\label{eq:isofactor}
\end{align}
In the orthonormal traceless operator coordinates, $J$ has eigenvalues
\begin{equation}\label{eq:isospectrum}
 v\quad(\hbox{multiplicity }5),\qquad
 2u-v\quad(\hbox{multiplicity }3),\qquad
 \det J=v^5(2u-v)^3.
\end{equation}
\end{proposition}
\begin{proof}
The two eigenspaces in~\eqref{eq:isospectrum} are, after complexification, the symmetric-traceless matrices and antisymmetric matrices in the Cartesian basis. Substitution in~\eqref{eq:cartesian}, or contraction against an orthonormal basis of these spaces, gives the stated eigenvalues.

For transparency, the polynomial identity in~\eqref{eq:isofactor} has a small exact computational proof. Write $H=P-Q_0$, $K=Q_0$; the scalar $-vI/3$ does not affect the commutators. Let $H_{12},H_{23},K_{12},K_{23}$ denote their embeddings and set
\[
 C_{20}=[H_{12},H_{23}],\quad
 C_{11}=[H_{12},K_{23}]+[K_{12},H_{23}],\quad
 C_{02}=[K_{12},K_{23}].
\]
Expanding $D$ in the monomials $u^3,u^2v,uv^2,v^3$ gives, respectively,
\begin{align*}
 &[H_{12}+H_{23},C_{20}],\\
 &[H_{12}+H_{23},C_{11}]+[K_{12}+K_{23},C_{20}],\\
 &[H_{12}+H_{23},C_{02}]+[K_{12}+K_{23},C_{11}],\\
 &[K_{12}+K_{23},C_{02}].
\end{align*}
Applying $9\pi_1\pi_3$ produces the four integer matrices
$9W,0,-9W,0$. The independent integer checker verifies each entry of these four matrices; the symbolic checker independently verifies the full two-variable identity. The matrix $W$ is given exactly in the evidence. In zero-based row-major indexing, $W_{0,4}=1/3$, so it is nonzero. Its $96$ nonzero entries comprise $48$ of absolute value $1/3$, $24$ of absolute value $2/3$ and $24$ of absolute value $1$. Thus
\[
 \|W\|_F^2=48/9+24(4/9)+24=40.
\]
Finally, $\Ob$ is rotationally invariant and odd under endpoint exchange. The two-spin decomposition $V_1\tens V_1=V_0\oplus V_1\oplus V_2$ is multiplicity-free, so any invariant two-site operator commutes with exchange. Hence $\Ob=0$. Its four cubic coefficient matrices are also checked to vanish exactly.
\end{proof}

\begin{lemma}[Both rank-deficient lines are leading-rigid]\label{lem:rankdef}
The interactions $h(1,0)$ and $h(1,2)$ satisfy Definition~\ref{def:rigid}. Consequently every nonzero member of~\eqref{eq:blbq} is leading-rigid.
\end{lemma}
\begin{proof}
Equation~\eqref{eq:isospectrum} leaves only $v=0$ and $v=2u$ outside full rank. On the first line the active local factors span the antisymmetric matrices, whose action on $\C^3$ is irreducible. On the second they span all symmetric-traceless matrices. Either set has only scalar matrices in its commutant, proving the boundary injectivity. The exact checker also verifies both one-sided maps have rank $8$.

For the pair condition, use the ordered basis
\begin{equation}\label{eq:realbasis}
 E_{01},E_{02},E_{10},E_{12},E_{20},E_{21},
 \diag(1,-1,0),\diag(0,1,-1)
\end{equation}
for $\mathfrak{sl}_3(\C)$. The map
\[
 \mathcal L_h(a,b)=[a_{12},h_{23}]-[h_{12},b_{23}]
\]
has $128$ columns and $729$ rows in row-major matrix coordinates. The supplied certificate indices select $127\times127$ minors with the following exact residues:
\begin{center}
\begin{tabular}{lrr}
\toprule
Integer interaction & Minor order & Determinant modulo $101$\\
\midrule
$P-Q_0=h(1,0)$ &127&73\\
$3P+3Q_0-2I=3h(1,2)$ &127&59\\
\bottomrule
\end{tabular}
\end{center}
Both minors are nonzero, so the rank is at least $127$ over $\Q$ and $\C$. The nonzero domain vector $(h,h)$ lies in the kernel, so the rank is at most $127$. The kernel is therefore exactly its span, proving the pair condition. Both checkers reconstruct the maps and their minors rather than trusting a recorded rank. Nonzero scaling preserves these properties, and Lemma~\ref{lem:fullrank} covers the rest of the family.
\end{proof}

\begin{theorem}[Arbitrary-field exclusion and its sharp interaction boundary]\label{thm:blbq}
Fix $(u,v)\ne(0,0)$. If
\begin{equation}\label{eq:nonintegrable}
 u(u^2-v^2)\ne0,
\end{equation}
then every periodic $N\ge6$ and every choice of Hermitian fields $f_1,\ldots,f_N$ has no conserved quantity of exact range $3\le k\le\lfloor N/2\rfloor$.

If $u(u^2-v^2)=0$, the zero-field chain has the nonzero conserved range-three quantity $C_N$. Thus, for this nonzero interaction family, \emph{some choice of fields admits a conserved quantity in the stated higher-range interval if and only if}
\[
 u=0\quad\hbox{or}\quad u=v\quad\hbox{or}\quad u=-v.
\]
The existence statement holds for every $N\ge6$ with zero fields; it does not assert preservation under arbitrary fields on those lines.
\end{theorem}
\begin{proof}
For~\eqref{eq:nonintegrable}, combine Lemma~\ref{lem:rankdef}, Proposition~\ref{prop:isocubic} and Theorem~\ref{thm:symmetry}.

On the three remaining lines $\Delta=0$. We show why this implies conservation of $C_N$ here, although no such general implication was used in Theorem~\ref{thm:quotient}. Since $D$ has no component with both endpoints traceless, it is a sum of neighbouring two-site and one-site terms. Rotational invariance removes all traceless one-site terms. Its two-site components are invariant and hence swap-even. Reversal oddness then forces $D=X_{23}-X_{12}$ for a two-site $X$. The translation sum telescopes to zero. Together with the four-site Jacobi cancellation, this proves $[H_N^{(2)},C_N]=0$. Leading rigidity and Proposition~\ref{prop:reduction} give $C_N\ne0$ of exact range three. The verifier additionally checks the full telescoping identity at the representatives $(0,1),(1,1),(1,-1)$; homogeneity covers the lines.
\end{proof}

\subsection{All remaining short-range charges}

\begin{theorem}[Exact local commutant in the excluded family]\label{thm:isocommutant}
Under~\eqref{eq:nonintegrable}, let
\[
 \mathfrak s_f=\{x\in\spanC\{S_x,S_y,S_z\}:[x,f_i]=0\text{ for every }i\}.
\]
Then
\begin{equation}\label{eq:isocommutant}
 \{Q:[H_N,Q]=0,\ \len(Q)\le\lfloor N/2\rfloor\}
 =\spanC\{I,H_N\}+\left\{\sum_i x_i:x\in\mathfrak s_f\right\},
\end{equation}
where $x_i$ is the copy of the \emph{same} local matrix $x$ at every site.
\end{theorem}
\begin{proof}
Theorem~\ref{thm:blbq} reduces a charge to range at most two, and leading rigidity fixes its two-site part to $cH_N^{(2)}$. After subtracting $cH_N$ and a scalar, the charge is $\sum_i x_i$ for traceless matrices $x_i$, initially allowed to differ. Its two-site commutator equations are
\begin{equation}\label{eq:bond}
 [x_i\tens I+I\tens x_{i+1},h]=0.
\end{equation}
Define $D_x$ by $[x,F_b]=\sum_a(D_x)_{ab}F_a$. The bilinear trace pairing gives $D_x^T=-D_x$, also for complex $x$. Equation~\eqref{eq:bond} and its transpose imply
\begin{equation}\label{eq:J2}
 D_{x_i}J-JD_{x_{i+1}}=0,
 \qquad [D_{x_i},J^2]=0.
\end{equation}
This implication does not require $J$ to be invertible.

In the Cartesian basis, the two eigenvalues of $J^2$ differ by
$(2u-v)^2-v^2=4u(u-v)\ne0$. Hence $D_{x_i}$ preserves the antisymmetric and symmetric-traceless matrix spaces separately. Decompose $x_i$ by transpose parity. Its antisymmetric part preserves those spaces under commutation, whereas its symmetric part exchanges them. The latter must therefore commute with all antisymmetric matrices; their scalar commutant forces this traceless symmetric part to vanish. Thus every $x_i$ is a spin generator.

For a spin generator, $D_{x_i}$ commutes with $J$. The first equation in~\eqref{eq:J2} becomes $(D_{x_i}-D_{x_{i+1}})J=0$. At least one of the two interaction sectors is nonzero. The spin action on either active sector has trivial kernel: equivalently, a spin generator commuting with all antisymmetric matrices or with all symmetric-traceless matrices is zero. Thus $x_i=x_{i+1}$ around the ring. The remaining one-site equations are exactly $[x,f_i]=0$ separately at each site. Conversely those conditions make the displayed uniform spin sum commute with both interactions and fields, proving equality.
\end{proof}

At zero fields this leaves the three components of total spin. Fields with no common commuting spin direction leave only $I$ and $H_N$. This is a finite exact condition on the fields, not an assumption about chaotic spectra or thermalisation.

\subsection{AKLT as a sharp check on the projection}
For $u=1$, $v=1/3$,
\[
 \det J=\frac{125}{6561}\ne0,\qquad
 \Ob(h)=0,\qquad
 \|\Delta(h)\|_F^2=40(1-1/9)^2=\frac{2560}{81}>0.
\]
Thus the earlier gap certificate is genuinely blind at full rank. The new symmetry-evaluated quotient detects this interaction and excludes higher-range local charges for \emph{arbitrary} on-site potentials. This strengthens the field scope of the established AKLT nonintegrability result; it does not claim the zero-field result as new.

\section{The cubic quick test and unrestricted genericity}\label{sec:quick}

\begin{proposition}[The gap projection is a quotient obstruction]\label{prop:gap}
For any symmetric $h$, $\Tr_2T_h(a,f)=0$. Therefore $\Ob(h)\ne0$ implies $\mathcal U(h)\ne0$. If $J$ is invertible, this excludes all higher-range charges for arbitrary fields.
\end{proposition}
\begin{proof}
Overlapping two-site terms have a commutator at the middle site, so their middle trace is zero. The density $C$ has three traceless factors; commutation with an on-site field retains tracelessness at every site. Thus tracing the middle factor kills both terms of $T_h$. Since $\Tr_2\Delta=\Ob$, membership $\Delta\in\im T_h$ would imply $\Ob=0$.
\end{proof}

The actual middle-scalar projection inserts $I/d$ between the two factors of $\Ob$. This factor is required because the trace is unnormalised.

\subsection{Coefficient formula and Hermiticity}
Write $[F_a,F_b]=i\sum_c f_{abc}F_c$, with real totally antisymmetric structure constants. Define the symmetric quadratic matrix
\begin{equation}\label{eq:K}
 K(J)_{rc}=\sum_{a,e,b,t} f_{rae}f_{cbt}J_{ab}J_{et}.
\end{equation}
Then direct middle-trace contraction gives
\begin{equation}\label{eq:cubiccoeff}
 \Ob(h)=\sum_{r,s}[J,K(J)]_{rs}F_r\tens F_s.
\end{equation}
Indeed, the two terms of $D$ contribute $-KJ$ and $JK$, respectively; symmetry of $J$ makes $K$ symmetric. Its commutator coefficient matrix is real and skew-symmetric. This means that the \emph{tensor operator} is Hermitian and exchange-odd, not anti-Hermitian: every tensor factor $F_r\tens F_s$ is Hermitian.

For $d=2$, $f_{abc}=\sqrt2\epsilon_{abc}$ and $K(J)=4\operatorname{adj}(J)$. Thus $\Ob=0$ for every symmetric qubit interaction. The general quotient remains defined, but no nonvanishing qubit gap obstruction is asserted.

\subsection{Residual short charges beyond simple spectrum}

\begin{lemma}[Abelian interaction eigenspaces]\label{lem:abelian}
Suppose the higher-range charges have been excluded and $h$ is leading-rigid. If every eigenspace of $J^2$, regarded as a subspace of traceless local matrices, is abelian under commutation, then
\begin{equation}\label{eq:IH}
 \{Q:[H_N,Q]=0,\ \len(Q)\le\lfloor N/2\rfloor\}
 =\spanC\{I,H_N\}
\end{equation}
for every choice of fields. Simple spectrum of $J^2$ is a sufficient special case.
\end{lemma}
\begin{proof}
As in Theorem~\ref{thm:isocommutant}, it remains to solve~\eqref{eq:bond}. Equation~\eqref{eq:J2} makes $D_{x_i}$ block diagonal in the eigenspaces of $J^2$. For $u,v$ in the same such eigenspace, trace invariance gives
$\Tr(v[x_i,u])=\Tr(x_i[u,v])=0$. Abelianity therefore makes every block zero. Hence $D_{x_i}=0$, and the centreless property of $\g$ gives $x_i=0$.
\end{proof}

\begin{theorem}[Unrestricted generic nonintegrability]\label{thm:generic}
For every $d\ge3$, an open dense full-Lebesgue-measure set of real symmetric traceless interactions has local commutant~\eqref{eq:IH} for every periodic $N\ge6$ and all site-dependent Hermitian fields. For $d=3$, one explicit sufficient nonzero polynomial is
\begin{equation}\label{eq:poly}
 P(J)=\det J\;\Disc_t\det(tI-J^2)\;\Tr(\Ob(h)^2),
 \qquad \deg P=8+112+6=126.
\end{equation}
\end{theorem}
\begin{proof}
Appendix~\ref{app:old} gives an exact qutrit witness with $\det J\ne0$, simple spectrum of $J^2$, and $\|\Ob\|_F^2=81\,962\,496$. Thus $P$ is not the zero polynomial, and its nonzero set has the stated topological and measure properties. Proposition~\ref{prop:gap} and Lemma~\ref{lem:abelian} give the conclusion. The same interaction set works for all chain lengths and fields; there is no size-dependent exceptional set to intersect. Its product with all uniform fields gives the 44-parameter statement.

For $d>3$, embed the eight normalised qutrit matrices in a $3\times3$ block and complete them to an orthonormal traceless Hermitian basis. The embedded obstruction is unchanged and nonzero. Give the complementary coordinates distinct, sufficiently small positive diagonal coupling coefficients. The interaction matrix becomes invertible with positive distinct eigenvalues and the obstruction remains nonzero. The analogous polynomial is therefore nonzero in every such dimension.
\end{proof}

This generic result is inherited from the version-0.1 argument, now situated inside the stronger quotient. It does not imply coverage of every symmetry-constrained subfamily. The isotropic blindness of its last factor is precisely why Sections~\ref{sec:symmetry}--\ref{sec:blbq} are necessary.

\section{Generic exclusion on both blind diagonal hypersurfaces}\label{sec:hypersurfaces}

Use the unnormalised Hermitian qutrit basis
\begin{align}\label{eq:G}
 &(G_1,G_2,G_3,G_4,G_5,G_6,G_7,G_8)\\[-2mm]
 &\quad=(X_{01},Y_{01},\diag(1,-1,0),X_{02},Y_{02},X_{12},Y_{12},\diag(1,1,-2)),\nonumber
\end{align}
where $X_{ab}=E_{ab}+E_{ba}$ and $Y_{ab}=-iE_{ab}+iE_{ba}$. The Gram diagonal is $(2,2,2,2,2,2,2,6)$. For
$h(j)=\sum_{a=1}^8 j_aG_a\tens G_a$, the exact old identity is
\begin{equation}\label{eq:diaggap}
 \Ob(h(j))=4(j_3-3j_8)(j_4j_5-j_6j_7)
       (G_3\tens G_8-G_8\tens G_3).
\end{equation}
The complete eight-variable identity is verified symbolically in the bundle. Its zero set is the union of
\[
 Z_C=\{j_3=3j_8\},\qquad Z_P=\{j_4j_5=j_6j_7\}.
\]
The new quotient is not identically zero on either component.

\begin{theorem}[Two certified blind-locus families]\label{thm:blind}
Each of $Z_C$ and $Z_P$ contains a nonempty relative Zariski-open set of interactions for which~\eqref{eq:IH} holds for every $N\ge6$ and all on-site fields. These sets are dense and of full measure in the respective regular real loci after exclusion of the additional algebraic degeneracies described in the proof.
\end{theorem}
\begin{proof}
The two exact witnesses are
\begin{equation}\label{eq:blindwitness}
 j^C=(1,2,24,4,5,6,7,8),\qquad
 j^P=(1,2,3,4,15,5,12,8).
\end{equation}
Every entry is nonzero, so both interactions have full rank. Both have $\Ob=0$. In the $1458$-row real-coordinate convention of Section~\ref{sec:quotient}, form the augmented matrix $[9T_h\mid9\Delta]$. The $28$ correction columns use $G_a\tens G_b-G_b\tens G_a$, ordered lexicographically with $a<b$, followed by $8$ fields $G_a$, then $\Delta$. The supplied $37\times37$ minors have determinants
\begin{center}
\begin{tabular}{lrrr}
\toprule
Witness & $\rank_{\Q}T_h$ & $\rank_{\Q}[T_h\mid\Delta]$ & Minor modulo $101$\\
\midrule
$j^C$ &36&37&43\\
$j^P$ &36&37&64\\
\bottomrule
\end{tabular}
\end{center}
Since $T_h$ has at most $36$ columns, either nonzero augmented minor already proves $\Delta\notin\im T_h$ over $\R$; the exact rational ranks provide a corroborating check. By Theorem~\ref{thm:quotient}, higher-range charges are excluded.

Each fixed minor is a homogeneous polynomial of degree
$28\cdot1+8\cdot2+3=47$ in $j$. The displayed witnesses show that its restriction to the corresponding hypersurface is nonzero. The hyperplane $Z_C$ and the determinant hypersurface $Z_P$ are irreducible, and their regular real loci are Zariski dense. Thus nonvanishing defines the stated relatively generic sets.

For $Z_P$, exclude vanishing coordinates and coincidences between the absolute values of $2j_1,\ldots,2j_7,6j_8$. These conditions hold at $j^P$, so they remove only proper algebraic subsets. The simple-spectrum case of Lemma~\ref{lem:abelian} applies. On $Z_C$, the forced equality $2j_3=6j_8$ leaves a two-dimensional eigenspace spanned by the commuting diagonal matrices $G_3,G_8$. Excluding every other absolute-value coincidence leaves all remaining eigenspaces one-dimensional. These conditions hold at $j^C$, and the general form of Lemma~\ref{lem:abelian} applies. The exact local bond map has rank $16$ at both witnesses, an additional check of the absence of residual on-site generators.
\end{proof}

This statement concerns the two hypersurfaces inside the fixed eight-parameter diagonal-$G$ family. It is not a decomposition of the full 36-parameter zero set of $\Ob$, and it does not classify their intersections or every exceptional point.

\section{Algorithms, benchmarks and evidence}\label{sec:evidence}

The public program \texttt{certify\_field\_uniform.py} accepts an exact rational symmetric $8\times8$ coefficient matrix in the $G$ basis. It first checks full interaction rank, then uses the gap obstruction when available. If the gap vanishes, it compares the exact ranks of $T_h$ and $[T_h\mid\Delta]$. A separate $16$-column local bond map determines whether any one-site generators remain. For the bilinear--biquadratic family, a named input mode applies Theorems~\ref{thm:blbq} and~\ref{thm:isocommutant}, including the rank-deficient cases.

All ranks used to issue a result are exact rational ranks. The stored nonzero minors use the prime $101$ only as \emph{lower-rank certificates} over characteristic zero; no vanishing modular minor is interpreted as a characteristic-zero rank upper bound. Input floats, asymmetric coefficient arrays, invalid rationals and non-finite values are rejected. The older \texttt{certify\_interaction.py} remains as a compatibility implementation of the quick test.

The matrix sizes depend on local dimension, not chain length. The formal input size still includes the bit lengths of rational coefficients, so no claim of constant arithmetic or bit complexity is made. The file \texttt{BENCHMARKS.md} records actual wall-clock timings and the environment for this release. In the supplied historical run, the seven original suites passed in $24.626$ seconds, including process startup. The publication revision adds an eighth suite for portable witnesses and tamper rejection; fresh timings are retained in the machine-readable publication receipt. These are illustrative measurements, not cross-machine guarantees; the machine-readable report records the environment and individual timings.

The original bundle includes seven runnable suites: the inherited integer and symbolic checks, the inherited certifier controls, the new symbolic/modular extension checks, an independent integer implementation of the isotropic calculation and rank certificates, dense structure-constant regressions, and controls for the stronger public certifier. The new independent implementation imports none of the other submitted modules. This is implementation diversity within the present work, not an independent human review.

\subsection{Portable quotient witnesses}
The optional \texttt{--witness} output records the input and its digest, the exact coordinate convention, and either a sparse rational dual functional $y$ with $y^T(9T_h)=0$ and $y^T(9\Delta)=1$, or rational coefficients $x$ with $(9T_h)x+9\Delta=0$. The separate \texttt{portable\_witness.py} command reconstructs the matrices from a separately supplied intended interaction and checks these identities by multiplication, without rank computation. It shares the matrix-construction code, so it is not an independent implementation. The ordinary certifier output is labelled an exact-recomputation report. Neither portable identity certifies leading rigidity, the residual commutant or the imported all-range theorem; consistency certifies only almost-conservation. Tamper controls alter coefficients, dimensions, coordinate conventions and input binding. Assertion-based legacy suites refuse Python optimisation rather than silently dropping tests.

\subsection{What the certificates establish}
The two $127$-minors and their known kernel vectors establish exact pair-kernel dimensions. The $37$-minors certify two new positive obstruction witnesses and the nonzero polynomial restrictions used for genericity. The full two-variable identity and four separately constructed homogeneous coefficient matrices prove the isotropic factorisation for all real parameters, rather than sampling it. Three additional dense rational interactions check the direct nested commutator against the structure-constant formula. The isotropic family and three qubit examples are negative controls for the \emph{old} obstruction; the three integrable isotropic lines are positive controls for a conserved current.

The all-range theorem, symmetry averaging, quotient equivalence and residual-commutant arguments are mathematical proofs in this manuscript. They are not reduced to numerical experiments. Exact finite calculations support the stated algebraic steps, but the proof as a whole has not been formalised in a proof assistant.

\section{Limitations and release boundary}\label{sec:limits}

The complete statement in Theorem~\ref{thm:blbq} classifies whether \emph{there exists a field choice} admitting a higher-range charge for an isotropic interaction. It does not classify all individual field configurations on the three charge-admitting lines. The general quotient similarly decides existence of an almost-conserved three-site candidate, not exact integrability when its obstruction vanishes. The unrestricted rank-deficient qutrit class, the full quotient-zero set, and the complete exceptional-set classification remain outside the proved results.

The range bound is $3\le k\le\lfloor N/2\rfloor$ on periodic chains. We do not exclude charges of larger contiguous range, nonlocal spectral projectors, polynomials in global charges, quasilocal tails, effective integrability in a projected low-energy sector, or approximate conservation at large fields. Nor do we infer thermalisation, the eigenstate thermalisation hypothesis, or level-spacing statistics.

Version 0.1 received the supplied favourable referee report. The present version addresses its scholarly revisions and adds theorems that require new specialist scrutiny; that earlier report must not be presented as refereeing these additions. The targeted literature search found no direct prior statement of the combined arbitrary-field theorem and quotient. This is not a historical-priority certificate. The publication revision responds to the supplied version-1 review without changing the main theorems. The review is not authenticated external peer review, and its separately mentioned sandbox audit archive was not supplied for this publication replay. The new portable-witness checks are producer-side evidence. Archival identifiers and exact payload identity are given in the accompanying citation and release records.

\appendix
\section{Self-contained pair-rigidity proof}\label{app:pair}

\begin{lemma}[Pair rigidity for $\mathfrak{sl}_d$]\label{lem:pair}
If $L,R\in\End(\g)$ satisfy $[Lx,y]=[x,Ry]$ for all $x,y\in\g$, then $L=R=c\id$.
\end{lemma}
\begin{proof}
Use the complex bilinear trace form, not a conjugate transpose. Trace invariance and nondegeneracy give $L^*[y,z]=[Ry,z]$. Interchanging $y,z$ also gives $L^*[y,z]=[y,Rz]$. Thus
\[
 [Rx,y]=[x,Ry],\qquad [Rx,x]=0.
\]
For a regular traceless diagonal $D$, the second relation says $RD$ is diagonal; by linearity and density this holds for every traceless diagonal $D$. With $i\ne j$, apply the first relation to $D,E_{ij}$ to obtain
\[
 RE_{ij}=c_{ij}E_{ij}+D_{ij},\qquad
 (RD)_i-(RD)_j=c_{ij}(D_i-D_j),
\]
where $D_{ij}$ is traceless diagonal. Reversing a root gives $c_{ij}=c_{ji}$. For three distinct indices, independence of $D_i-D_j$ and $D_j-D_k$ gives $c_{ik}=c_{ij}=c_{jk}$. Thus all coefficients equal a single $c$ for $d\ge3$; the reversed-root identity suffices for $d=2$. Hence $RD=cD$ on diagonal matrices.

Apply the relation to $E_{ij},E_{kl}$. The scalar terms cancel, leaving
\[
 ((D_{ij})_k-(D_{ij})_l)E_{kl}
 =((D_{kl})_j-(D_{kl})_i)E_{ij}.
\]
Distinct ordered pairs make the two matrix units independent. For fixed $(i,j)$ this removes all diagonal differences except possibly its own; the reversed ordered pair removes that difference as well. Hence $D_{ij}$ is scalar and traceless, and is zero. We have $R=c\id$. The original relation then makes $(L-c\id)x$ central for every $x$, so $L=c\id$.
\end{proof}

This proof is included for inspectability, with the quasicentroid attribution in Section~\ref{sec:rigidity}. The inherited supplement also verifies a distinct $512\times128$ formulation of this abstract pair lemma for $d=3$: a $127$-minor has determinant $32$ modulo $101$, and $(\id,\id)$ is the known kernel vector. That certificate is separate from the two interaction-specific pair certificates in Lemma~\ref{lem:rankdef}.

\section{Retained cubic witness and a certified perturbation ball}\label{app:old}

For~\eqref{eq:G} and $j^*=(1,2,3,4,5,6,7,8)$, the interaction is
\begin{equation}\label{eq:hstar}
 h_*=\begin{pmatrix}
11&0&0&0&-1&0&0&0&-1\\
0&5&0&3&0&0&0&0&0\\
0&0&-16&0&0&0&9&0&0\\
0&3&0&5&0&0&0&0&0\\
-1&0&0&0&11&0&0&0&-1\\
0&0&0&0&0&-16&0&13&0\\
0&0&9&0&0&0&-16&0&0\\
0&0&0&0&0&13&0&-16&0\\
-1&0&0&0&-1&0&0&0&32
\end{pmatrix}.
\end{equation}
Its orthonormal coefficient matrix is $J_* =\diag(2,4,6,8,10,12,14,48)$. Direct integer arithmetic gives
\begin{align*}
 \Ob(h_*)&=1848(G_3\tens G_8-G_8\tens G_3)\\
 &=3696\diag(0,1,-1,-1,0,1,1,-1,0),\\
 \|\Ob(h_*)\|_F^2&=6(3696)^2=81\,962\,496.
\end{align*}
The factorised formula~\eqref{eq:diaggap} follows by expanding
\[
 \sum_{a,c,e}j_aj_cj_e\Tr(G_e[G_a,G_c])
 \bigl([G_e,G_a]\tens G_c+G_a\tens[G_e,G_c]\bigr).
\]
Only the $(3,8)$ and $(8,3)$ basis coefficients survive. The symbolic verifier checks all coefficients as polynomials in eight indeterminates.

\begin{proposition}[A field-uniform robust neighbourhood]\label{prop:ball}
If $h$ is Hermitian, swap-symmetric, has zero one-site partial traces and $\|h-h_*\|_F\le1/10$, then its local commutant is~\eqref{eq:IH} for every choice of fields and every $N\ge6$.
\end{proposition}
\begin{proof}
Orthonormality gives $\|J-J_*\|_F=\|h-h_*\|_F$. The eigenvalues of $J_*$ are positive with minimum $2$ and minimum gap $2$, so the stated perturbation preserves positivity and simplicity, and hence simplicity of $J^2$.

For interactions $a,b$, set $M_a=\|a\|_{\op}$, $M_b=\|b\|_{\op}$ and $\delta=\|a-b\|_F$. Expanding the difference of the nested commutators and using
$\|[X,Y]\|_F\le2\|X\|_{\op}\|Y\|_F$,
$\|X\tens I\|_F=\sqrt d\|X\|_F$ and
$\|\Tr_2X\|_F\le\sqrt d\|X\|_F$ gives
\[
 \|\Ob(a)-\Ob(b)\|_F
 \le8d\delta(M_a^2+M_aM_b+M_b^2).
\]
For clarity, writing $A=a_{12}$, $B=a_{23}$ and primed letters for $b$, the difference is
$[A+B-A'-B',[A,B]]+[A'+B',[A,B]-[A',B']]$.
The two bounds before tracing are $8\sqrt d\delta M_a^2$ and $8\sqrt d\delta M_b(M_a+M_b)$, respectively.

The maximum absolute row sum of~\eqref{eq:hstar} is $34$, so $M_{h_*}\le34$ and $M_h\le35$. Consequently
\[
 \|\Ob(h)-\Ob(h_*)\|_F
 \le\frac{24}{10}(35^2+35\cdot34+34^2)=8570.4
 <\sqrt{81\,962\,496}.
\]
The gap obstruction remains nonzero; Proposition~\ref{prop:gap} and Lemma~\ref{lem:abelian} finish the proof.
\end{proof}

\section{Certificate conventions and reproducibility}\label{app:certs}

All matrix indices in the evidence are zero-based. Tensor bases order the left site first and flatten matrices row by row. The isotropic certificate \texttt{isotropic\_W.json} stores exact rational strings. The two interaction-pair certificates store the integer two-site matrix, the ordered basis~\eqref{eq:realbasis}, the selected rows and columns, and the nonzero determinant residue. The $127$-minors use the first $64$ columns for $a$ and the next $64$ for $b$ in~\eqref{eq:pairlocal}.

The two blind-locus files store the eight diagonal $G$ coefficients, $37$ selected rows, all $37$ augmented columns, and the determinant modulo $101$. Their matrix is $[9T_h\mid9\Delta]$ in the real-then-imaginary flattening of Section~\ref{sec:quotient}. The factors of $9$ clear the endpoint projection denominators. Both rank-deficient pair matrices and both augmented minors are rebuilt from their inputs when checked.

From the bundle directory, install the pinned symbolic dependency and run:
\begin{quote}\ttfamily
python -m pip install -r requirements.txt\\
python verify.py\\
python certify\_field\_uniform.py evidence/AKLT\_input.json
\end{quote}
The independent integer checker and the inherited integer checker need only the Python standard library. Verification reports record actual results, environment and timings. The checksum manifest protects the release contents; verification can regenerate timing-bearing reports without implying their old hashes must remain unchanged. The source audit distinguishes inspected theorem statements and proofs from metadata-only or abstract-only comparisons. Original software is MIT-licensed; original prose and data are dedicated under CC0-1.0, with third-party exceptions in the licence map. Archival citation metadata uses the standing scholarly creator Anonymous and the reserved version DOI printed on the title page; publication itself requires separate provider readback.

\section*{Acknowledgement of method and review status}
This manuscript was developed with AI-assisted mathematical reasoning and exact computation. The supplied version-0.1 report motivated the symmetry analysis and literature corrections. It is included for provenance; its assessment does not extend automatically to the new results. Independent review of the quotient equivalence, the arbitrary-field symmetry step and the complete spin-1 statement is the principal remaining assurance step.

\begin{thebibliography}{99}
\bibitem{HokkyoTest}
Hokkyo, A. (2025). \emph{Rigorous test for quantum integrability and nonintegrability} (Version 3, 13 April 2025) [Preprint]. arXiv:2501.18400.\newline
\url{https://arxiv.org/abs/2501.18400v3}

\bibitem{ParkLee}
Park, H. K., \& Lee, S. (2024). \emph{Proof of nonintegrability of the spin-1 bilinear-biquadratic chain model} [Preprint]. arXiv:2410.23286.\newline
\url{https://arxiv.org/abs/2410.23286}

\bibitem{HSC}
Hokkyo, A., Yamaguchi, M., \& Chiba, Y. (2025). Absence of nontrivial local conserved quantities in the spin-1 bilinear-biquadratic chain and its anisotropic extensions. \emph{Physical Review Research, 7}, 043297.\newline
\url{https://doi.org/10.1103/vjtq-sz21}

\bibitem{SW}
Sengoku, S., \& Watanabe, H. (2026). \emph{Absence of nontrivial local conserved quantities in a class of $U(1)$-symmetric spin-1 chains} [Preprint]. arXiv:2608.17548.\newline
\url{https://arxiv.org/abs/2608.17548}

\bibitem{SY}
Shiraishi, N., \& Yamaguchi, M. (2026). Dichotomy theorem separating complete integrability and nonintegrability of isotropic spin chains. \emph{Physical Review B, 113}, L241111.\newline
\url{https://doi.org/10.1103/lcn4-tbcm}

\bibitem{HokkyoLaw}
Hokkyo, A. (2026). Integrability from a single conservation law in quantum spin chains. \emph{Physical Review B, 114}, 055108.\newline
\url{https://doi.org/10.1103/y66t-t8zw}

\bibitem{SSI}
Sanatani, M., Shiraishi, N., \& Ishii, F. (2026). \emph{A simple necessary and sufficient condition for Yang--Baxter integrability} [Preprint]. arXiv:2607.29660.\newline
\url{https://arxiv.org/abs/2607.29660v1}

\bibitem{LegerLuks}
Leger, G. F., \& Luks, E. M. (2000). Generalized derivations of Lie algebras. \emph{Journal of Algebra, 228}, 165--203. The quasicentroid statements are consulted through~\cite{ABR}, p.~8.\newline
\url{https://www.sciencedirect.com/science/article/pii/S0021869399982509}

\bibitem{ABR}
Andruszkiewicz, R. R., Brzezi\'nski, T., \& Radziszewski, K. (2024). \emph{Lie affgebras vis-\`a-vis Lie algebras} [Preprint]. arXiv:2409.01996.\newline
\url{https://arxiv.org/abs/2409.01996}

\bibitem{GrabowskiMathieu}
Grabowski, M. P., \& Mathieu, P. (1995). Integrability test for spin chains. \emph{Journal of Physics A: Mathematical and General, 28}, 4777--4798.\newline
\url{https://doi.org/10.1088/0305-4470/28/17/013}
\end{thebibliography}
\end{document}
